% Options for packages loaded elsewhere
% Options for packages loaded elsewhere
\PassOptionsToPackage{unicode}{hyperref}
\PassOptionsToPackage{hyphens}{url}
\PassOptionsToPackage{dvipsnames,svgnames,x11names}{xcolor}
%
\documentclass[
  11pt,
  letterpaper,
  DIV=11,
  numbers=noendperiod]{scrartcl}
\usepackage{xcolor}
\usepackage[top=2.5cm,left=2.5cm,right=2.5cm,heightrounded]{geometry}
\usepackage{amsmath,amssymb}
\setcounter{secnumdepth}{4}
\usepackage{iftex}
\ifPDFTeX
  \usepackage[T1]{fontenc}
  \usepackage[utf8]{inputenc}
  \usepackage{textcomp} % provide euro and other symbols
\else % if luatex or xetex
  \usepackage{unicode-math} % this also loads fontspec
  \defaultfontfeatures{Scale=MatchLowercase}
  \defaultfontfeatures[\rmfamily]{Ligatures=TeX,Scale=1}
\fi
\usepackage{lmodern}
\ifPDFTeX\else
  % xetex/luatex font selection
  \setmainfont[]{Latin Modern Roman}
  \setsansfont[]{Latin Modern Roman}
\fi
% Use upquote if available, for straight quotes in verbatim environments
\IfFileExists{upquote.sty}{\usepackage{upquote}}{}
\IfFileExists{microtype.sty}{% use microtype if available
  \usepackage[]{microtype}
  \UseMicrotypeSet[protrusion]{basicmath} % disable protrusion for tt fonts
}{}
\usepackage{setspace}
\makeatletter
\@ifundefined{KOMAClassName}{% if non-KOMA class
  \IfFileExists{parskip.sty}{%
    \usepackage{parskip}
  }{% else
    \setlength{\parindent}{0pt}
    \setlength{\parskip}{6pt plus 2pt minus 1pt}}
}{% if KOMA class
  \KOMAoptions{parskip=half}}
\makeatother
% Make \paragraph and \subparagraph free-standing
\makeatletter
\ifx\paragraph\undefined\else
  \let\oldparagraph\paragraph
  \renewcommand{\paragraph}{
    \@ifstar
      \xxxParagraphStar
      \xxxParagraphNoStar
  }
  \newcommand{\xxxParagraphStar}[1]{\oldparagraph*{#1}\mbox{}}
  \newcommand{\xxxParagraphNoStar}[1]{\oldparagraph{#1}\mbox{}}
\fi
\ifx\subparagraph\undefined\else
  \let\oldsubparagraph\subparagraph
  \renewcommand{\subparagraph}{
    \@ifstar
      \xxxSubParagraphStar
      \xxxSubParagraphNoStar
  }
  \newcommand{\xxxSubParagraphStar}[1]{\oldsubparagraph*{#1}\mbox{}}
  \newcommand{\xxxSubParagraphNoStar}[1]{\oldsubparagraph{#1}\mbox{}}
\fi
\makeatother


\usepackage{longtable,booktabs,array}
\usepackage{calc} % for calculating minipage widths
% Correct order of tables after \paragraph or \subparagraph
\usepackage{etoolbox}
\makeatletter
\patchcmd\longtable{\par}{\if@noskipsec\mbox{}\fi\par}{}{}
\makeatother
% Allow footnotes in longtable head/foot
\IfFileExists{footnotehyper.sty}{\usepackage{footnotehyper}}{\usepackage{footnote}}
\makesavenoteenv{longtable}
\usepackage{graphicx}
\makeatletter
\newsavebox\pandoc@box
\newcommand*\pandocbounded[1]{% scales image to fit in text height/width
  \sbox\pandoc@box{#1}%
  \Gscale@div\@tempa{\textheight}{\dimexpr\ht\pandoc@box+\dp\pandoc@box\relax}%
  \Gscale@div\@tempb{\linewidth}{\wd\pandoc@box}%
  \ifdim\@tempb\p@<\@tempa\p@\let\@tempa\@tempb\fi% select the smaller of both
  \ifdim\@tempa\p@<\p@\scalebox{\@tempa}{\usebox\pandoc@box}%
  \else\usebox{\pandoc@box}%
  \fi%
}
% Set default figure placement to htbp
\def\fps@figure{htbp}
\makeatother


% definitions for citeproc citations
\NewDocumentCommand\citeproctext{}{}
\NewDocumentCommand\citeproc{mm}{%
  \begingroup\def\citeproctext{#2}\cite{#1}\endgroup}
\makeatletter
 % allow citations to break across lines
 \let\@cite@ofmt\@firstofone
 % avoid brackets around text for \cite:
 \def\@biblabel#1{}
 \def\@cite#1#2{{#1\if@tempswa , #2\fi}}
\makeatother
\newlength{\cslhangindent}
\setlength{\cslhangindent}{1.5em}
\newlength{\csllabelwidth}
\setlength{\csllabelwidth}{3em}
\newenvironment{CSLReferences}[2] % #1 hanging-indent, #2 entry-spacing
 {\begin{list}{}{%
  \setlength{\itemindent}{0pt}
  \setlength{\leftmargin}{0pt}
  \setlength{\parsep}{0pt}
  % turn on hanging indent if param 1 is 1
  \ifodd #1
   \setlength{\leftmargin}{\cslhangindent}
   \setlength{\itemindent}{-1\cslhangindent}
  \fi
  % set entry spacing
  \setlength{\itemsep}{#2\baselineskip}}}
 {\end{list}}
\usepackage{calc}
\newcommand{\CSLBlock}[1]{\hfill\break\parbox[t]{\linewidth}{\strut\ignorespaces#1\strut}}
\newcommand{\CSLLeftMargin}[1]{\parbox[t]{\csllabelwidth}{\strut#1\strut}}
\newcommand{\CSLRightInline}[1]{\parbox[t]{\linewidth - \csllabelwidth}{\strut#1\strut}}
\newcommand{\CSLIndent}[1]{\hspace{\cslhangindent}#1}



\setlength{\emergencystretch}{3em} % prevent overfull lines

\providecommand{\tightlist}{%
  \setlength{\itemsep}{0pt}\setlength{\parskip}{0pt}}



 


\usepackage{setspace}
\linespread{1.2}
\usepackage{changepage}
\usepackage{longtable}
\usepackage{nameref}
\usepackage{hyperref}
\usepackage{booktabs}
\usepackage{tabularray}
\usepackage{float}
\KOMAoption{captions}{tableheading}
\makeatletter
\@ifpackageloaded{caption}{}{\usepackage{caption}}
\AtBeginDocument{%
\ifdefined\contentsname
  \renewcommand*\contentsname{Table of contents}
\else
  \newcommand\contentsname{Table of contents}
\fi
\ifdefined\listfigurename
  \renewcommand*\listfigurename{List of Figures}
\else
  \newcommand\listfigurename{List of Figures}
\fi
\ifdefined\listtablename
  \renewcommand*\listtablename{List of Tables}
\else
  \newcommand\listtablename{List of Tables}
\fi
\ifdefined\figurename
  \renewcommand*\figurename{Figure}
\else
  \newcommand\figurename{Figure}
\fi
\ifdefined\tablename
  \renewcommand*\tablename{Table}
\else
  \newcommand\tablename{Table}
\fi
}
\@ifpackageloaded{float}{}{\usepackage{float}}
\floatstyle{ruled}
\@ifundefined{c@chapter}{\newfloat{codelisting}{h}{lop}}{\newfloat{codelisting}{h}{lop}[chapter]}
\floatname{codelisting}{Listing}
\newcommand*\listoflistings{\listof{codelisting}{List of Listings}}
\makeatother
\makeatletter
\makeatother
\makeatletter
\@ifpackageloaded{caption}{}{\usepackage{caption}}
\@ifpackageloaded{subcaption}{}{\usepackage{subcaption}}
\makeatother
\usepackage{bookmark}
\IfFileExists{xurl.sty}{\usepackage{xurl}}{} % add URL line breaks if available
\urlstyle{same}
\hypersetup{
  pdftitle={Populist Pressure, Public Opinion, and Central Bank Independence},
  colorlinks=true,
  linkcolor={black},
  filecolor={Maroon},
  citecolor={black},
  urlcolor={black},
  pdfcreator={LaTeX via pandoc}}


\title{Populist Pressure, Public Opinion, and Central Bank Independence}
\usepackage{etoolbox}
\makeatletter
\providecommand{\subtitle}[1]{% add subtitle to \maketitle
  \apptocmd{\@title}{\par {\large #1 \par}}{}{}
}
\makeatother
\subtitle{A Theoretical Framework on How Public Opinion Mediates
Populist Attacks on Central Bank Independence}
\author{Nikolaos Vichos\textsuperscript{}}
\date{December 16, 2025}
\begin{document}
\maketitle
\begin{abstract}
\begin{spacing}{0.95}
Populist leaders have increasingly targeted independent regulatory agencies, yet existing research has focused primarily on the institutional, legal, and market-based constraints shaping the effectiveness of such attacks. This paper shifts attention to the role of public opinion. Focusing on central banks as a most-likely case of an independent regulatory agency, it develops a theoretical framework that conceptualizes populist attacks as public-facing bargaining processes in which de facto independence is contingent on mass politics. Building on work that models conflicts between governments and central banks as episodes of strategic escalation, the paper argues that the success of populist pressure depends on two dimensions of public opinion: support for the incumbent launching the attack and trust in the central bank. When public opinion aligns with the government and is skeptical of technocratic authority, populist pressure is more likely to translate into monetary accommodation; when it does not, central banks are better positioned to resist. By integrating insights from the literatures on regulatory politicization, populism, and judicial delegitimation, the paper highlights public opinion as a central but underexplored mediator of populist challenges to independent institutions.
\end{spacing}
\end{abstract}


\setstretch{1.19}
\pagenumbering{gobble}

\newpage{}

\thispagestyle{empty}

\tableofcontents

\newpage{}

\pagenumbering{arabic}

\section{\texorpdfstring{\textbf{Introduction}}{Introduction}}\label{introduction}

Independent regulatory agencies (IRAs) are seemingly under constant
attack: whether one looks at Trump pressuring the Federal Reserve to
lower interest rates (\citeproc{ref-mcferran2025}{McFerran and
McNicholas 2025}), Bolsonaro curtailing the independence of Brazil's
environmental agencies (\citeproc{ref-bersch2024}{Bersch and Lotta
2024}), or the PiS's vocal criticism of Polish Ombudsman Adam Bodnar
(\citeproc{ref-mazur2021}{Mazur 2021}), the overarching picture is one
of IRAs' independence being continuously threatened. Most importantly,
the culprit tends to be the same in most cases: populist leaders of the
executive branch.

This trend is an important cause for concern. Even though independent
regulators should not be insulated from criticism on the grounds of
their often limited accountability, transparency, and responsibility
(\citeproc{ref-maggetti2010}{Maggetti 2010};
\citeproc{ref-tripathi2018}{Tripathi 2018}), their independence also
enables them to effectively execute their role. Non-independent
politicized regulators tend to suffer systematic declines in
organizational capacity, as political interference reshapes the
incentives, composition, and behavior of the individuals operating
within them (\citeproc{ref-moynihan2025}{D. P. Moynihan 2025}).
Prioritizing political loyalty over expertise leads to the recruitment
of more politically connected but less competent personnel, while
politicized work environments increase turnover and induce the exit of
experienced officials (\citeproc{ref-colonnelli2020}{Colonnelli, Prem,
and Teso 2020}; \citeproc{ref-richardson2019}{Richardson 2019};
\citeproc{ref-xu2018}{Xu 2018}). High turnover also erodes institutional
memory and managerial competence, as new political appointees typically
lack task-specific experience and remain in office for short periods
(\citeproc{ref-gallo2012}{Gallo and Lewis 2012}). This weakened
bureaucratic capacity that results from a lack of independence also
increases vulnerability to waste and politically motivated spending, as
reduced oversight and reliance on politically aligned contractors raise
costs and facilitate favoritism (\citeproc{ref-dahlstrom2021}{Dahlström,
Fazekas, and Lewis 2021}). Finally, beyond capacity, democratic
responsiveness is also weakened, as politicized agencies are not only
less responsive to information requests from the public or legislatures
(\citeproc{ref-wood2017}{Wood and Lewis 2017}) but also less able to
achieve policy goals (\citeproc{ref-moynihan2025}{D. P. Moynihan 2025}).

Despite its concerning nature, however, this trend might also seem
unsurprising when closely considering the specific nature of both the
`perpetrators' (the populist executive) and the `victim' (IRAs). As
non-elected bodies that are staffed with technical experts and tasked
with overseeing essential tenets of public policy yet remain formally
independent from political control, IRAs can easily become the
caricature of the ``elites'' in populists' Manichean worldview of a
society divided ``into two homogeneous and antagonistic groups, `the
pure people' versus `the corrupt elite'\,''
(\citeproc{ref-mudde2004}{Mudde 2004, 543}). This specific framing has
been used to target a broad range of different IRAs, including central
banks, environmental agencies, or competition regulators
(\citeproc{ref-bernatt2025}{Bernatt 2025};
\citeproc{ref-bersch2024}{Bersch and Lotta 2024};
\citeproc{ref-jung2025}{Jung 2025})

Ironically, one aspect that remains unclear is the role of the
public---the primary `audience' of the populists---in enabling or
constraining the effectiveness of these attacks. By occupying a central
role in the populists' narrative defending the `real' people's interests
by launching these attacks, the public could serve as a decisive
mechanism that mediates the success or failure of efforts to erode
institutional autonomy. Despite existing literature having examined how
third factors like institutions, fiscal rules, or international
financial markets might constrain or exacerbate the impact of populists'
attacks on independent regulators, the role of public opinion has been
notably overlooked.

To understand how public opinion might mediate the effectiveness of
populist attacks on IRAs, I focus specifically on central banks. Though
they might not be the most `typical' case of a regulator under attack
due to their more well-known role among the public, I view central banks
as worth studying exactly due to their prominence as the most visible
and widely recognized IRA (\citeproc{ref-mcferran2025}{McFerran and
McNicholas 2025}). The high salience of their mandate---which includes
the promotion of price stability and, in some cases, full employment not
only means that attempts at political influence might be more likely
(\citeproc{ref-benoit2021}{Benoît 2021};
\citeproc{ref-ringquist2003}{Ringquist, Worsham, and Eisner 2003}) but
also makes it more likely that the public will pay attention during
these attacks. This makes them a `most likely' case study where the
mechanism behind the mediating role of the public is most likely to be
observed (\citeproc{ref-flyvbjerg2006}{Flyvbjerg 2006}). I thus pose the
following research question:

\begin{adjustwidth}{1em}{1em}
\emph{Does public opinion mediate the effectiveness of populist attacks against the independence of central banks?}
\end{adjustwidth}

To address this question, I propose a theoretical framework that
conceptualizes populist attacks on central banks as public-facing
bargaining processes in which de facto independence is contingent on
mass politics. Building on existing work that models conflicts between
governments and central banks as episodes of strategic escalation
(\citeproc{ref-gavin2023}{Gavin and Manger 2023}), I argue that the
effectiveness of populist pressure depends not only on institutional and
economic constraints, but also on public opinion. In particular, I
theorize that public support for the incumbent launching the attack and
public trust in the central bank jointly condition whether populist
pressure translates into monetary accommodation.

\section{\texorpdfstring{\textbf{Literature
Review}}{Literature Review}}\label{literature-review}

This paper draws from existing literature on (i) the politicization of
regulatory processes, (ii) populism and the relationship between
populist leaders and bureaucracies, and (iii) specific literature on
attacks against central banks and the factors that constrain or
exacerbate their frequency or intensity. After a brief discussion of
some key concepts and definitions, I present each of the three
aforementioned strands.

\subsection{Definitions}\label{definitions}

To understand what a politicized IRA is, one must also understand what
independence means within the context of regulatory agencies.

An agency's independence refers to ``the degree to which the agency
takes day-to-day decisions without the interference of politicians---in
terms of the offering of inducements or threats---and/or the
consideration of political preferences''
(\citeproc{ref-hanretty2013}{Hanretty and Koop 2013, 196}). This
independence can be categorized as formal/de jure, meaning that the
legal framework governing the agency endows it with a certain degree of
independence from politics, or informal/de facto, which refers to the
extent to which ``the agency operates independently from politics in
practice'' (\citeproc{ref-hanretty2013}{Hanretty and Koop 2013, 196}).

This definition, along with its underlying dichotomy, applies to central
banks as well. Most central banks are formally independent, enjoying
important legal autonomy to ``adjust interest rates and other policy
tools to fight inflation without political interference''
(\citeproc{ref-gavin2023}{Gavin and Manger 2023, 1190}). Their informal
independence, however, often comes under strain
(\citeproc{ref-goodhart2018}{Goodhart and Lastra 2018};
\citeproc{ref-jung2025}{Jung 2025}). This paper thus focuses on central
banks' de facto independence: this is because of the specific nature of
the strategies populists employ when undermining the independent
institutions, maintaining their legal/formal structure while severely
undercutting their capacity to deliver and resist
(\citeproc{ref-bauer2021}{Bauer et al. 2021};
\citeproc{ref-yesilkagit2025}{Yesilkagit, Papadopoulos, and Coban 2025};
\citeproc{ref-moynihan2021}{D. Moynihan 2021}).

\subsection{Politicization of Regulatory
Processes}\label{politicization-of-regulatory-processes}

For the purposes of this paper, the opposite of a \emph{de facto}
independent regulatory agency is a politicized one. I refrain from
adopting overly narrow definitions that focus exclusively on the
staffing of agencies by politicized bureaucrats\footnote{For example,
  Ennser-Jedenastik (\citeproc{ref-ennser-jedenastik2016}{2016, 509})
  defines politicization as the appointment, retention, promotion, or
  dismissal (if possible) of bureaucrats based on political criteria
  rather than merit.''} or overly broad definitions that risk
`overdetecting' politicization by focusing on phenomena as broad as the
``increasing awareness of, and mobilization around, regulation by
non-majoritarian institutions'' (\citeproc{ref-onoda2024}{Onoda 2024,
1348}). Instead, I propose a definition of politicization that
implicitly recognizes politics as an integral and ever-present part of
regulatory processes (\citeproc{ref-short2021}{Short 2021}) (thereby
avoiding the collapse of routine political oversight into politicization
itself) but refrains from overly narrow criteria. Politicization thus
refers to a\emph{ny systematic influence on regulatory decision-making
that prioritizes political considerations over expertise or established
rules.} This can take place through changes in not only the staffing of
agencies, but also their structure, resource allocation, or the norms
governing them (\citeproc{ref-moynihan2021}{D. Moynihan 2021}).

The study of regulatory politicization has followed distinct
trajectories across geographic contexts, particularly between the United
States and Western Europe. As Benoît (\citeproc{ref-benoit2021}{2021})
argues, U.S.-based scholarship has tended to emphasize politicization by
relying on principal--agent models that foreground political control and
agency loss, often portraying regulators as highly responsive to
political principals. However, this literature may overstate the extent
of politicization by focusing on politicians' incentives while
underestimating their limited capacity to monitor and steer regulatory
action, as well as regulators' own agency and reputational constraints.
In contrast, European scholarship has often underemphasized
politicization, focusing instead on agencies' legal independence,
expertise, and internal dynamics, and frequently treating political
intervention as marginal or ineffective. This divide points toward the
need for a broader understanding of politicization that moves beyond
formal control and recognizes how regulators respond to political
signals from multiple actors, without necessarily sacrificing their
autonomy.

Politicization and informal independence are mutually exclusive, at
least for the purposes of this paper; however, the relationship of
politicization with \emph{formal} independence is significantly more
complex. At first sight, greater levels of formal independence should
lead to lower politicization. This is supported by the findings of
Benoît and Szilagyi (\citeproc{ref-benoit2023}{2023}): in studying
ex-ante and ex-post legislative involvement in the creation of 48 IRAs
in France, they find that agencies that were formally designed to allow
greater participation by legislators were indeed subject to greater
levels of political scrutiny. Nevertheless, findings by
Ennser-Jedenastik (\citeproc{ref-ennser-jedenastik2016}{2016}) and Onoda
(\citeproc{ref-onoda2024}{2024}) partly contradict this. The former, in
a large-N study looking at the relationship between measures of formal
independence and the partisan background of bureaucrats, finds that high
levels of legal independence lead to an increased likelihood of
individuals with party ties being appointed. This suggests that
institutions might create incentives to undermine non-interference
norms. Onoda's (2024) findings represent a middle solution connecting
the two aforementioned authors: focusing on drug rationing policies in
England and France, he finds that the degree of initial
independence/delegation can determine politicization-driven policy
change. Whereas low levels of initial delegation lead to greater policy
stability because elected politicians intervened to prevent unpopular
decisions, highly delegated settings suffered from important policy
instability due to the counter-mobilization against the unpopular
measures taken by the regulator.

Beyond formal independence, the extent of politicization also depends on
more public-facing factors such as the regulator's reputation or the
salience and technical complexity of the issue at hand. In studying
IRAs' responses to public opinion, Maor, Gilad, and Bloom
(\citeproc{ref-maor2013}{2013}) find that the Israeli banking regulator
kept silent on issues (i) where it enjoys a stronger reputation or (ii)
that fall outside its jurisdiction but responded to opinions targeting
issues where it was reputationally weaker or that touched upon its core
operations. Similarly, Ringquist, Worsham, and Eisner
(\citeproc{ref-ringquist2003}{2003}) examine four different regulatory
areas that vary in complexity and salience and find that public salience
has a positive effect on the willingness of legislators to intervene to
direct the behavior of IRAs, especially in less complex areas like
consumer protection. These findings suggest that, despite the seemingly
technical and detached nature of regulation, public opinion remains
central, with not only IRAs but also their principals (politicians)
paying close attention to changes in public mood.

The centrality of public opinion is likely to increase even more as a
result of the ``declining appeal of valence politics''
(\citeproc{ref-hay2019}{Hay and Benoît 2019, 389}). Focusing on the 2016
Brexit referendum, Hay and Benoît (\citeproc{ref-hay2019}{2019}) argue
that Remain's framing of Brexit as a ``valence issue''---technical and
presumably uncontroversial questions around which there should be
widespread agreement---proved less effective at mobilizing electoral
participation than Leave's positional framing, which cast Brexit as a
matter of political choice and popular sovereignty. More broadly, their
analysis situates this dynamic within a longer-term erosion of ``expert
paternalism,'' (\citeproc{ref-hay2020}{Hay 2020}) in which trust in
insulated expertise as a basis for depoliticized governance has
weakened, and policy domains once treated as technical have become sites
of political contestation. It is within this context of
re-politicization that populist actors are able to mobilize public
opinion against expert-led institutions, a theme taken up explicitly in
the populism literature.

\subsection{Populism and Attacks on Bureaucracies and Regulatory
Independence}\label{populism-and-attacks-on-bureaucracies-and-regulatory-independence}

Rather than conceiving of populism as a political strategy
(\citeproc{ref-weyland2017}{Weyland 2017}) or a style of economic
policy-making (\citeproc{ref-sachs1989}{Sachs 1989}), I instead rely on
the ideational approach defining populism as ``an ideology that
considers society to be ultimately separated into two homogeneous and
antagonistic groups, `the pure people' versus `the corrupt elite', and
which argues that politics should be an expression of the volonté
générale (general will) of the people'' (\citeproc{ref-mudde2004}{Mudde
2004}). Crucially, this definition goes against authors---including in
the domain of regulation---who assume a specific set of policies such as
economic expansionism Goodhart and Lastra
(\citeproc{ref-goodhart2018}{2018}). Instead, I recognize that the
governing style and policy substance of populist governments are likely
to be quite diverse, especially between populists on opposite sides of
the spectrum (\citeproc{ref-bugaric2019}{Bugaric 2019};
\citeproc{ref-mudde2013}{Mudde and Kaltwasser 2013}).

At the same time, however, I also argue that their Manichean framing and
anti-elitist/anti-pluralist bent also predispose them to adopt a
particularly hostile position toward counter-majoritarian institutions.
This is most commonly applied to the study of liberal democratic
backsliding, with populists rejecting its `liberal' components such as
the rule of law, minority protections, or the separation of powers;
however, I argue that this applies equally well to IRAs not only due to
their counter-majoritarian character but also as a result of their high
levels of expertise, making them easy targets of rhetoric framing them
as `out-of-touch technocrats' (\citeproc{ref-mudde2021}{Mudde 2021}). In
that sense, I concur with Mudde (\citeproc{ref-mudde2021}{2021}) in that
populism can be viewed as an ``illiberal democratic response to
undemocratic liberalism'' (592). Populists thus seek to re-politicize
certain issues that had been taken away from the public and into the
hands of technically skilled regulators. As a result, within the context
of this study, I view populism's effects on the independence of
regulators as \emph{politicization on steroids}: driven by a Manichean
worldview, populists who seek to re-politicize issues delegated to IRAs
portray regulators as corrupt, out-of-touch elites acting against the
will of the people. This moralized, zero-sum understanding of political
conflict lowers the threshold for resorting to rhetoric and
interventions that directly undermine regulatory agencies' autonomy.

\subsubsection{Strategies of Populists in
Power}\label{strategies-of-populists-in-power}

Literature in this area typically examines how populists in power seek
to undermine or reshape bureaucratic autonomy once electoral victory has
been secured. A useful organizing framework is provided by Bauer et al.
(\citeproc{ref-bauer2021}{2021}), who argue that illiberal populist
governments typically pursue one of three broad strategies toward the
administrative state: sideline it, ignore it, or use it.

\begin{enumerate}
\def\labelenumi{\arabic{enumi}.}
\tightlist
\item
  Populists may seek to \textbf{sideline} bureaucracies by constructing
  alternative decision-making structures closer to the executive,
  surrounding themselves with loyal experts, or creating parallel
  institutions that bypass existing agencies altogether
  (\citeproc{ref-bauer2021}{Bauer et al. 2021};
  \citeproc{ref-mazur2021}{Mazur 2021}).
\item
  A second strategy consists of ignoring the bureaucracy. Here,
  populists refrain from direct confrontation and instead govern in an
  ad hoc and personalized manner, relying on a narrow circle of
  ideologically aligned advisors while allowing bureaucratic routines to
  continue largely unchecked. While less visibly confrontational, such
  neglect can still weaken bureaucratic influence by excluding agencies
  from key decision-making arenas and information flows
  (\citeproc{ref-bauer2021}{Bauer et al. 2021}).
\item
  Finally, populists may seek to use the bureaucracy instrumentally.
  This strategy involves efforts to reshape bureaucratic incentives and
  cultures, encouraging compliance through politicized appointments,
  informal pressure, or selective enforcement. In its most extreme form,
  instrumentalization may support a broader authoritarian project by
  aligning administrative power with executive preferences
  (\citeproc{ref-bauer2021}{Bauer et al. 2021};
  \citeproc{ref-moynihan2021}{D. Moynihan 2021}).
\end{enumerate}

Crucially, Bauer et al. (\citeproc{ref-bauer2021}{2021}) show that these
strategies operate along several interrelated dimensions. Populists may
centralize administrative structures, redistribute resources to weaken
disfavored units, politicize personnel through purges or patronage,
undermine professional norms of neutrality, and hollow out
accountability by reducing transparency and oversight. Comparative
evidence from Poland illustrates how these dimensions can be combined in
practice: the PiS government simultaneously centralized civil service
management, altered recruitment standards, reallocated budgets, and
violated norms of political neutrality, leading to a measurable decline
in administrative quality (\citeproc{ref-mazur2021}{Mazur 2021}).

Country cases beyond Poland further illustrate these strategies in
action. In the U.S., the Trump administration launched an unprecedented
campaign against IRAs by summarily removing agency heads on political
grounds, stacking commissions with loyalists, or leaving key positions
vacant (\citeproc{ref-mcferran2025}{McFerran and McNicholas 2025}). It
expanded White House oversight over regulatory agendas by subjecting
IRAs to ex-ante and ex-post review by the Office of Management and
Budget, weakening their traditionally independent rulemaking and
enforcement authority. Finally, the Federal Reserve is repeatedly
targeted through public attacks, threats of dismissal, and the firing of
a member of its Board of Governors, signaling that even the most
insulated agencies were not beyond political reach. Similar dynamics
emerged in Brazil, where Bolsonaro gradually transformed environmental
governance by targeting individual bureaucrats, narrowing arenas of
contestation, and inducing silence through formal and informal
sanctions, ultimately repurposing institutional safeguards to
consolidate executive control (\citeproc{ref-bersch2024}{Bersch and
Lotta 2024}).

Importantly, the effects of populist strategies are not confined to the
internal functioning of bureaucracies. Recent work by Dobbins and
Labanino (\citeproc{ref-dobbins2025}{2025}) highlights how populist
attacks can reshape agencies' external environments, including their
relationships with civil society, markets, and international networks:
in post-communist contexts, illiberal rule alters patterns of
consultation between regulators and interest groups, often privileging
politically aligned actors while excluding neutral or oppositional
organizations.

While much of this literature focuses on bureaucracies and civil
services, IRAs occupy a more ambiguous position. As Yesilkagit,
Papadopoulos, and Coban (\citeproc{ref-yesilkagit2025}{2025}) argue,
IRAs are often less susceptible to direct hierarchical control due to
their formal insulation, yet this very insulation renders them
vulnerable to targeted delegitimization.

\subsubsection{Populists vs the Judiciary: Delegitimization
Strategies}\label{populists-vs-the-judiciary-delegitimization-strategies}

This logic closely parallels findings from the literature on populist
attacks against the judiciary. Like IRAs, courts are formally insulated
from direct political control in order to preserve impartiality, yet
this insulation simultaneously exposes them to strategies relying less
on immediate institutional capture and more on public-facing efforts to
erode trust and delegitimize expert authority.

Köker et al. (\citeproc{ref-kuxf6ker2025}{2025}) shows that this
vulnerability allows populists to weaken judicial independence even
without holding executive power. By systematically criticizing court
decisions, individual judges, and judicial institutions as a whole,
populists deploy verbal delegitimization strategies that frame the
judiciary as a corrupt and self-serving elite obstructing the will of
the people. These attacks are not primarily aimed at reforming judicial
behavior, but at undermining the normative foundations of judicial
authority by eroding public trust. Given courts' limited enforcement
capacity, sustained rhetorical assaults can substantially weaken their
effectiveness in practice, even in the absence of formal constitutional
or legal change.

Whether such delegitimization translates into meaningful institutional
erosion, however, depends critically on how it is received by the
public. Magalhães and Garoupa (\citeproc{ref-magalhuxe3es2024}{2024})
finds that assaults on judicial independence in Hungary, Poland, and
Turkey reduce trust primarily among citizens ideologically distant from
governing parties, while supporters of populist leaders remain largely
unaffected. Political alignment thus conditions the impact of populist
attacks, as individuals are more willing to tolerate or discount
institutional erosion when it is carried out by leaders they perceive as
legitimate representatives of the people. Public opinion thus acts as a
mediating mechanism, filtering the effectiveness of verbal attacks on
independent institutions; this dynamic also motivates this paper's focus
on the public as an enabler or constraint of populist pressure.

\subsection{Central Banks and Factors Mediating Populist
Attacks}\label{central-banks-and-factors-mediating-populist-attacks}

Current work on how populists interact with central banks largely builds
on a broader political economy literature showing that politicians, on
average, favor lower interest rates due to electoral incentives,
distributional considerations, and concerns about debt servicing
(\citeproc{ref-ehrmann2011}{Ehrmann and Fratzscher 2011};
\citeproc{ref-lastra2001}{Lastra and Miller 2001}). However, rather than
altering statutes or mandates, politicians often rely on informal
pressure, public criticism, and attempts to influence expectations.

Gavin and Manger (\citeproc{ref-gavin2023}{2023}) show that populist
politicians are significantly more likely than non-populists to exert
public pressure on central banks, and that such pressure is often
effective in securing interest rate concessions. Importantly, these
conflicts rarely translate into irregular central bank governor
turnover, suggesting that de jure independence can coexist with
substantial de facto political influence. The authors also find that
populist leaders, especially on the right, are more willing to escalate
conflicts and incur economic or political costs, making brinkmanship a
more credible strategy in confrontations with central banks.

\subsection{Mediating factors}\label{mediating-factors}

Most importantly for the purposes of this paper, existing research has
emphasized that the effectiveness of populist pressure on central banks
is conditional on a range of institutional and economic constraints.
Financial markets are a key moderating factor: when they react
negatively to political interference, politicians are less likely to
escalate pressure, and central banks are more likely to resist
(\citeproc{ref-gavin2023}{Gavin and Manger 2023}). Fiscal institutions
constitute another important constraint, though counterintuitively: Jung
(\citeproc{ref-jung2025}{2025}) finds that under stricter fiscal rules,
populist leaders are more likely to undermine central bank independence
as a substitute for constrained fiscal policy, particularly when
short-term economic stimulus is politically attractive.

Beyond markets and fiscal rules, legal and institutional safeguards like
mandate clarity, appointment rules, and removal protection also shape
the effectiveness of populist interference in central banking
(\citeproc{ref-goodhart2018}{Goodhart and Lastra 2018}). However, this
form remains an incomplete shield: central banks' authority rests not
only on legal insulation, but also on broader societal acceptance of
their role, goals, and expertise. As Goodhart and Lastra
(\citeproc{ref-goodhart2018}{2018}) notes, challenges to central banks
often take the form of questioning their objectives, legitimacy, or
accountability, rather than direct statutory violations. Importantly,
populist leaders often choose to attack central banks openly and in
public, rather than relying solely on behind-the-scenes pressure or
institutional reform. In conjunction with the previously discussed
findings by Magalhães and Garoupa
(\citeproc{ref-magalhuxe3es2024}{2024}), Köker et al.
(\citeproc{ref-kuxf6ker2025}{2025}), and Maor, Gilad, and Bloom
(\citeproc{ref-maor2013}{2013}), this suggests that conflicts over
central bank independence are waged not only within institutional arenas
but also in the public sphere, pointing to public opinion as an
important yet underexplored mediating factor in the effectiveness of
populist attacks.

\section{Theoretical Framework}\label{theoretical-framework}

Drawing from Gavin and Manger (\citeproc{ref-gavin2023}{2023}), this
paper conceptualizes conflicts between populist governments and central
banks as public-facing episodes of crisis bargaining in which formally
independent central banks face pressure through public escalation rather
than legal reform. In these conflicts, de jure independence remains
intact, yet de facto independence becomes contingent on the strategic
interaction between politicians, central banks, and their relevant
audiences.

A key insight of Gavin and Manger (\citeproc{ref-gavin2023}{2023}) is
that populist politicians are more likely to initiate and escalate such
conflicts because they face lower political costs when doing so. Rather
than quietly negotiating or pursuing statutory change, populists
deliberately apply pressure in public, framing central banks as
technocratic elites acting against the will of the people. Whether this
strategy succeeds depends on the costs of escalation and resistance,
which existing work has primarily modeled through reactions by financial
markets.

This paper retains the logic of public pressure and brinkmanship,
shifting attention from markets to public opinion as the key
intermediary shaping conflict outcomes. Public pressure is not only an
instrument to coerce central banks but also a signal directed at voters,
showing that the populist leader is willing to clash with the `corrupt
elite.' As a result, audience reactions condition both politicians'
willingness to escalate and central banks' incentives to resist. Public
opinion thus plays a dual role: it imposes audience costs on politicians
who escalate unsuccessfully, and it affects the legitimacy costs central
banks incur when resisting pressure. Two dimensions of public opinion
are thus particularly salient.

\begin{enumerate}
\def\labelenumi{\arabic{enumi}.}
\tightlist
\item
  Support for the incumbent government: Public approval can shape the
  political costs of escalation. Populist leaders who enjoy strong
  public backing face lower audience costs when pressuring central banks
  and are therefore more willing to persist.
\item
  Public trust in the central bank: Because central banks rely on
  societal acceptance of their expertise and mandate
  (\citeproc{ref-goodhart2018}{Goodhart and Lastra 2018}), low public
  trust can weaken their capacity to withstand pressure without risking
  reputational damage.
\end{enumerate}

Together, these dynamics generate variation in when populist pressure
succeeds. When public opinion aligns with the incumbent and is skeptical
of the central bank, escalation becomes politically attractive and
resistance costly, increasing the likelihood of monetary accommodation.
When public support for the government is weak, or trust in the central
bank is high, public pressure is riskier and resistance more viable.
This leads to the following hypotheses:

\begin{enumerate}
\def\labelenumi{\arabic{enumi}.}
\item
  \emph{H1 (Populist Pressure Hypothesis):} Populist governments are
  more likely than non-populist governments to exert public pressure on
  central banks to ease monetary policy.
\item
  \emph{H2 (Public Opinion Conditioning Hypothesis):} The effect of
  populist pressure on central bank policy outcomes is conditioned by
  public opinion.
\item
  \emph{H2a (Government Support):} Populist pressure is more likely to
  result in interest rate cuts when public support for the incumbent
  government is high.
\item
  \emph{H2b (Central Bank Trust):} Populist pressure is more likely to
  result in interest rate cuts when public trust in the central bank is
  low.
\item
  \emph{H3 (Alignment Hypothesis):} The probability that a central bank
  accommodates populist pressure is highest when public opinion
  simultaneously favors the incumbent government and is skeptical of the
  central bank.
\end{enumerate}

\section{Conclusion}\label{conclusion}

This paper has developed a theoretical framework for understanding
populist attacks on central bank independence by placing public opinion
at the center of the analysis. Conceptualizing such attacks as
public-facing bargaining processes, it argues that de facto independence
depends not only on institutional design and economic constraints, but
also on mass politics. In particular, public support for the incumbent
and public trust in the central bank jointly condition whether populist
pressure translates into monetary accommodation.

While the paper is primarily theoretical, the framework lends itself to
empirical evaluation.\footnote{Initially, this formed part of the paper
  as a separate section. However, I removed it to remain consistent with
  the page limit.} At the cross-national level, its implications can be
examined quantitatively: for example, large-N analyses using
logit/probit models linking episodes of populist pressure to observable
monetary policy outcomes (decreasing vs maintaining/increasing interest
rates) under varying public opinion conditions. At the same time,
qualitative evidence, including expert interviews following
high-salience episodes such as the recent confrontations involving the
Federal Reserve and its repeated decisions between September 2025 and
December 2025 to lower interest rates (\citeproc{ref-shalal2025}{Shalal,
Jones, and Jones 2025}) can help illuminate the internal dynamics of
this mechanism. Together, these approaches offer a path for assessing
when public opinion constrains and when it enables populist challenges
to central bank independence.

\newpage{}

\section{Bibliography}\label{bibliography}

\phantomsection\label{refs}
\begin{CSLReferences}{1}{0}
\bibitem[\citeproctext]{ref-bauer2021}
Bauer, Michael W., B. Guy Peters, Jon Pierre, Kutsal Yesilkagit, and
Stefan Becker. 2021. {``Introduction: Populists, Democratic Backsliding,
and Public Administration.''} In, edited by B. Guy Peters, Jon Pierre,
Kutsal Yesilkagit, Michael W. Bauer, and Stefan Becker, 1--21.
Cambridge: Cambridge University Press.
\url{https://doi.org/10.1017/9781009023504.002}.

\bibitem[\citeproctext]{ref-benoit2021}
Benoît, Cyril. 2021. {``Politicians, Regulators, and Regulatory
Governance: The Neglected Sides of the Story.''} \emph{Regulation \&
Governance} 15 (S1): S8--22. \url{https://doi.org/10.1111/rego.12388}.

\bibitem[\citeproctext]{ref-benoit2023}
Benoît, Cyril, and Ana-Maria Szilagyi. 2023. {``Legislative Direction of
Regulatory Bureaucracies: Evidence from a Semi-Presidential System.''}
\emph{The Journal of Legislative Studies} 29 (2): 271--90.
\url{https://doi.org/10.1080/13572334.2021.1986281}.

\bibitem[\citeproctext]{ref-bernatt2025}
Bernatt, Maciej. 2025. {``Politicization of Competition Agencies: In
Search of an Analytical Framework Fit for Trump Era.''}
\url{https://doi.org/10.2139/ssrn.5375480}.

\bibitem[\citeproctext]{ref-bersch2024}
Bersch, Katherine, and Gabriela Lotta. 2024. {``Political Control and
Bureaucratic Resistance: The Case of Environmental Agencies in
Brazil.''} \emph{Latin American Politics and Society} 66 (1): 27--50.
\url{https://doi.org/10.1017/lap.2023.22}.

\bibitem[\citeproctext]{ref-bugaric2019}
Bugaric, Bojan. 2019. {``The Two Faces of Populism: Between
Authoritarian and Democratic Populism.''} \emph{German Law Journal} 20
(3): 390--400. \url{https://doi.org/10.1017/glj.2019.20}.

\bibitem[\citeproctext]{ref-colonnelli2020}
Colonnelli, Emanuele, Mounu Prem, and Edoardo Teso. 2020. {``Patronage
and Selection in Public Sector Organizations.''} \emph{American Economic
Review} 110 (10): 3071--99. \url{https://doi.org/10.1257/aer.20181491}.

\bibitem[\citeproctext]{ref-dahlstrom2021}
Dahlström, Carl, Mihály Fazekas, and David E. Lewis. 2021. {``Partisan
Procurement: Contracting with the United States Federal Government,
2003{\textendash}2015.''} \emph{American Journal of Political Science}
65 (3): 652--69. \url{https://doi.org/10.1111/ajps.12574}.

\bibitem[\citeproctext]{ref-dobbins2025}
Dobbins, Michael, and Rafael Labanino. 2025. {``From Agents of the
People to Agents of Authority? How Illiberal Populism Impacts
Interactions Between Regulatory Agencies and External Stakeholders.''}
\emph{Regulation \& Governance} 19 (3): 933--56.
\url{https://doi.org/10.1111/rego.12631}.

\bibitem[\citeproctext]{ref-ehrmann2011}
Ehrmann, Michael, and Marcel Fratzscher. 2011. {``Politics and Monetary
Policy.''} \emph{Review of Economics and Statistics} 93 (3): 941--60.
\url{https://doi.org/10.1162/REST_a_00113}.

\bibitem[\citeproctext]{ref-ennser-jedenastik2016}
Ennser-Jedenastik, Laurenz. 2016. {``The Politicization of Regulatory
Agencies: Between Partisan Influence and Formal Independence.''}
\emph{Journal of Public Administration Research and Theory} 26 (3):
507--18. \url{https://doi.org/10.1093/jopart/muv022}.

\bibitem[\citeproctext]{ref-flyvbjerg2006}
Flyvbjerg, Bent. 2006. {``Five Misunderstandings About Case-Study
Research.''} \emph{Qualitative Inquiry} 12 (2): 219--45.
\url{https://doi.org/10.1177/1077800405284363}.

\bibitem[\citeproctext]{ref-gallo2012}
Gallo, Nick, and David E. Lewis. 2012. {``The Consequences of
Presidential Patronage for Federal Agency Performance.''} \emph{Journal
of Public Administration Research and Theory} 22 (2): 219--43.
\url{https://doi.org/10.1093/jopart/mur010}.

\bibitem[\citeproctext]{ref-gavin2023}
Gavin, Michael, and Mark Manger. 2023. {``Populism and De Facto Central
Bank Independence.''} \emph{Comparative Political Studies} 56 (8):
1189--1223. \url{https://doi.org/10.1177/00104140221139513}.

\bibitem[\citeproctext]{ref-goodhart2018}
Goodhart, Charles, and Rosa Lastra. 2018. {``Populism and Central Bank
Independence.''} \emph{Open Economies Review} 29 (1): 49--68.
\url{https://doi.org/10.1007/s11079-017-9447-y}.

\bibitem[\citeproctext]{ref-hanretty2013}
Hanretty, Chris, and Christel Koop. 2013. {``Shall the Law Set Them
Free? The Formal and Actual Independence of Regulatory Agencies.''}
\emph{Regulation \& Governance} 7 (2): 195--214.
\url{https://doi.org/10.1111/j.1748-5991.2012.01156.x}.

\bibitem[\citeproctext]{ref-hay2020}
Hay, Colin. 2020. {``Brexistential Angst and the Paradoxes of Populism:
On the Contingency, Predictability and Intelligibility of Seismic
Shifts.''} \emph{Political Studies} 68 (1): 187--206.
\url{https://doi.org/10.1177/0032321719836356}.

\bibitem[\citeproctext]{ref-hay2019}
Hay, Colin, and Cyril Benoît. 2019. {``Brexit, Positional Populism, and
the Declining Appeal of Valence Politics.''} \emph{Critical Review} 31
(3-4): 389--404. \url{https://doi.org/10.1080/08913811.2019.1722531}.

\bibitem[\citeproctext]{ref-jung2025}
Jung, Hoyong. 2025. {``Central Bank Independence and Fiscal Rule Under
Populist Leader's Regime.''} \emph{European Journal of Political
Economy} 89 (September): 102728.
\url{https://doi.org/10.1016/j.ejpoleco.2025.102728}.

\bibitem[\citeproctext]{ref-kuxf6ker2025}
Köker, Philipp, Tilko Swalve, Merle Huber, Christoph Hönnige, and
Dominic Nyhuis. 2025. {``Populists Before Power: Delegitimization
Strategies Against Independent Judiciaries.''} \emph{Democratization} 32
(6): 1432--49. \url{https://doi.org/10.1080/13510347.2025.2470765}.

\bibitem[\citeproctext]{ref-lastra2001}
Lastra, Rosa, and Geoffrey Miller. 2001. {``Central Bank Independence in
Ordinary and Extraordinary Times.''} \emph{Central Bank Independence:
The Economic Foundations, the Constitutional Implications and Democratic
Accountability}, January.
\url{https://gretchen.law.nyu.edu/fac-chapt/2027}.

\bibitem[\citeproctext]{ref-magalhuxe3es2024}
Magalhães, Pedro C., and Nuno Garoupa. 2024. {``Populist Governments,
Judicial Independence, and Public Trust in the Courts.''} \emph{Journal
of European Public Policy} 31 (9): 2748--74.
\url{https://doi.org/10.1080/13501763.2023.2235386}.

\bibitem[\citeproctext]{ref-maggetti2010}
Maggetti, Martino. 2010. {``Legitimacy and Accountability of Independent
Regulatory Agencies: A Critical Review.''} \emph{Living Reviews in
Democracy}.

\bibitem[\citeproctext]{ref-maor2013}
Maor, Moshe, Sharon Gilad, and Pazit Ben-Nun Bloom. 2013.
{``Organizational Reputation, Regulatory Talk, and Strategic Silence.''}
\emph{Journal of Public Administration Research and Theory: J-PART} 23
(3): 581--608. \url{https://www.jstor.org/stable/24484861}.

\bibitem[\citeproctext]{ref-mazur2021}
Mazur, Stanisław. 2021. {``Public Administration in Poland in the Times
of Populist Drift.''} In, edited by B. Guy Peters, Jon Pierre, Kutsal
Yesilkagit, Michael W. Bauer, and Stefan Becker, 100--126. Cambridge:
Cambridge University Press.
\url{https://doi.org/10.1017/9781009023504.006}.

\bibitem[\citeproctext]{ref-mcferran2025}
McFerran, Lauren, and Celine McNicholas. 2025. {``Trump{'}s Assault on
Independent Agencies Endangers Us All.''} Washington, D.C.
\url{https://www.epi.org/publication/trumps-assault-on-independent-agencies-endangers-us-all/}.

\bibitem[\citeproctext]{ref-moynihan2021}
Moynihan, Donald. 2021. {``Populism and the Deep State: The Attack on
Public Service Under Trump.''} In, edited by B. Guy Peters, Jon Pierre,
Kutsal Yesilkagit, Michael W. Bauer, and Stefan Becker, 151--77.
Cambridge: Cambridge University Press.
\url{https://doi.org/10.1017/9781009023504.008}.

\bibitem[\citeproctext]{ref-moynihan2025}
Moynihan, Donald P. 2025. {``Rescuing State Capacity: Proceduralism, the
New Politicization, and Public Policy.''} \emph{Journal of Policy
Analysis and Management} 44 (2): 364--78.
\url{https://doi.org/10.1002/pam.22673}.

\bibitem[\citeproctext]{ref-mudde2004}
Mudde, Cas. 2004. {``The Populist Zeitgeist.''} \emph{Government and
Opposition} 39 (4): 541--63.
\url{https://doi.org/10.1111/j.1477-7053.2004.00135.x}.

\bibitem[\citeproctext]{ref-mudde2021}
---------. 2021. {``Populism in Europe: An Illiberal Democratic Response
to Undemocratic Liberalism (The Government and Opposition/Leonard
Schapiro Lecture 2019).''} \emph{Government and Opposition} 56 (4):
577--97. \url{https://doi.org/10.1017/gov.2021.15}.

\bibitem[\citeproctext]{ref-mudde2013}
Mudde, Cas, and Cristóbal Rovira Kaltwasser. 2013. {``Exclusionary Vs.
Inclusionary Populism: Comparing Contemporary Europe and Latin
America.''} \emph{Government and Opposition} 48 (2): 147--74.
\url{https://doi.org/10.1017/gov.2012.11}.

\bibitem[\citeproctext]{ref-onoda2024}
Onoda, Takuya. 2024. {``{`}Depoliticised{'} Regulators as a Source of
Politicisation: Rationing Drugs in England and France.''} \emph{Journal
of European Public Policy} 31 (5): 1346--67.
\url{https://doi.org/10.1080/13501763.2023.2176909}.

\bibitem[\citeproctext]{ref-richardson2019}
Richardson, Mark D. 2019. {``Politicization and Expertise: Exit, Effort,
and Investment.''} \emph{The Journal of Politics} 81 (3): 878--91.
\url{https://doi.org/10.1086/703072}.

\bibitem[\citeproctext]{ref-ringquist2003}
Ringquist, Evan J., Jeff Worsham, and Marc Allen Eisner. 2003.
{``Salience, Complexity, and the Legislative Direction of Regulatory
Bureaucracies.''} \emph{Journal of Public Administration Research and
Theory} 13 (2): 141--64. \url{https://doi.org/10.1093/jpart/mug013}.

\bibitem[\citeproctext]{ref-sachs1989}
Sachs, Jeffrey D. 1989. {``Social Conflict and Populist Policies in
Latin America.''} \url{https://doi.org/10.3386/w2897}.

\bibitem[\citeproctext]{ref-shalal2025}
Shalal, Andrea, Ryan Patrick Jones, and Ryan Patrick Jones. 2025.
{``Trump Thinks More Should Be Done to Lower Interest Rates, White House
Says.''} \emph{Reuters}, December.
\url{https://www.reuters.com/sustainability/sustainable-finance-reporting/trump-thinks-more-should-be-done-lower-interest-rates-white-house-says-2025-12-11/}.

\bibitem[\citeproctext]{ref-short2021}
Short, Jodi L. 2021. {``The Politics of Regulatory Enforcement and
Compliance: Theorizing and Operationalizing Political Influences.''}
\emph{Regulation \& Governance} 15 (3): 653--85.
\url{https://doi.org/10.1111/rego.12291}.

\bibitem[\citeproctext]{ref-tripathi2018}
Tripathi, Sudhanshu. 2018. {``Independent Regulatory Authorities:
Analysing Accountability, Responsibility and Transparency.''}
\emph{Indian Journal of Public Administration} 64 (3): 349--57.
\url{https://doi.org/10.1177/0019556118783098}.

\bibitem[\citeproctext]{ref-weyland2017}
Weyland, Kurt. 2017. {``Populism: A Political-Strategic Approach.''} In,
edited by Cristóbal Rovira Kaltwasser, Paul Taggart, Paulina Ochoa
Espejo, and Pierre Ostiguy, 0. Oxford University Press.
\url{https://doi.org/10.1093/oxfordhb/9780198803560.013.2}.

\bibitem[\citeproctext]{ref-wood2017}
Wood, Abby K, and David E Lewis. 2017. {``Agency Performance Challenges
and Agency Politicization.''} \emph{Journal of Public Administration
Research and Theory} 27 (4): 581--95.
\url{https://doi.org/10.1093/jopart/mux014}.

\bibitem[\citeproctext]{ref-xu2018}
Xu, Guo. 2018. {``The Costs of Patronage: Evidence from the British
Empire.''} \emph{American Economic Review} 108 (11): 3170--98.
\url{https://doi.org/10.1257/aer.20171339}.

\bibitem[\citeproctext]{ref-yesilkagit2025}
Yesilkagit, Kutsal, Ioannis Papadopoulos, and M. Kerem Coban. 2025.
{``The Global Retreat of the Regulatory State? The Populist Challenge
and the Future of Regulatory Governance.''}
\url{https://doi.org/10.31219/osf.io/wrhq4_v1}.

\end{CSLReferences}




\end{document}
